\documentclass[10pt,a4paper]{amsart}
\usepackage[margin=19mm]{geometry}
\usepackage{amsmath,amssymb,mathtools}
\usepackage[scr=rsfs,cal=euler]{mathalfa}
\usepackage{xfrac}
\usepackage[unicode,hidelinks,bookmarks=false]{hyperref}
\newcommand{\ie}{i.e.\ }
\newcommand{\cf}{cf.\ }
\newcommand{\resp}{respectively\ }
\newcommand{\Subsection}{\subsection}
\providecommand{\nobreakdash}{\nobreak\hbox{-}}
\setlength{\emergencystretch}{2em}
\multlinegap0pt
\usepackage{bm}
\usepackage{bbm}
\usepackage{dsfont}
\usepackage{xspace}
\usepackage{cancel}
\usepackage{mathtools}
\usepackage{amsthm}
\makeatletter
\@ifundefined{theorem}{\newtheorem{theorem}{Theorem}}{}
\@ifundefined{lemma}{\newtheorem{lemma}[theorem]{Lemma}}{}
\@ifundefined{proposition}{\newtheorem{proposition}[theorem]{Proposition}}{}
\makeatother
\def\R{\mathbb{R}}
\def\T{\mathcal{T}}
\def\X{\mathcal{X}}
\def\Y{\mathcal{Y}}
\def\aul{\overline{u}_l}
\def\aur{\overline{u}_r}
\def\avl{\overline{v}_l}
\def\avr{\overline{v}_r}
\def\ph{\mathbf{\Phi}}
\def\p{\partial}
\def\supt{\mathop{sup}\limits_{t \in [0,\mathcal{T}]}}
\def\ut{\tilde{u}}
\def\vsx{v^S_{\mathcal{X}}}
\def\vt{\tilde{v}}
\def\wh{\hat{w}}
\def\w{\mathbf{W}}
\newcommand{\norm}[1]{\left\lVert #1 \right\rVert}
\title[Stability of Large-Amplitude Viscous Shock Under Periodic Pe]{Stability of Large-Amplitude Viscous Shock Under Periodic Perturbation for 1-d Viscoelasticity with Non-Convex Constitutive Relations}
\author{Yu Mei, Peng Yuan}
\date{Source: arXiv:2507.21470v1 (CC BY 4.0)}
\begin{document}
\maketitle

\noindent\emph{Source and licence.} Stability of Large-Amplitude Viscous Shock Under Periodic Perturbation for 1-d Viscoelasticity with Non-Convex Constitutive Relations. Version arXiv:2507.21470v1 (CC BY 4.0), distributed under the Creative Commons Attribution 4.0 International licence, as recorded in the arXiv licence metadata for this exact version and shown on the abstract page https://arxiv.org/abs/2507.21470v1. Full rights notice: licenses/arxiv-2507.21470-CC-BY-4.0.txt.\par\smallskip

\noindent\textit{Editorial scope.} Counted unit: the Lemma whose source label is \texttt{construction} in the article (source0.tex lines 1122--1131, source label \texttt{construction}), delivered together with the article's own complete original proof of it (source0.tex lines 1132--1306, one \texttt{proof} environment of 175 lines). No mathematics is written, repaired or re-derived: statement and proof are the article's own text line for line. Required context retained uncounted: VNC-Ion (lines 171--176); initial (lines 196--198); cond-farfield (lines 200--206); 0-A (lines 238--240); decay (lines 327--342); XY (lines 408--415); XYdecay (lines 435--462); C1 (lines 453--455); new0mass (lines 495--503); mainresult (lines 531--549); system2 (lines 587--592); hterm (lines 594--596); newinitial (lines 607--609); mainpro (lines 612--624); small1 (lines 649--654); decay-error (lines 663--669); XYZ-cond (lines 704--706); transe (lines 733--741); Hbig (lines 965--967). Every definition, display and label of those lines is retained unchanged. Omitted from this package and not redistributed: 1--170, 177--195, 199--199, 207--237, 241--326, 343--407, 416--434, 456--494, 504--530, 550--586, 593--593, 597--606, 610--611, 625--648, 655--662, 670--703, 707--732, 742--964, 968--1121, 1307--1917. Nothing inside a retained excerpt is omitted or altered: every retained body is delivered exactly as the article's lines are written, so no omission receipt of any kind accompanies this package. Citations: the retained excerpts carry no citation command at all, so no bibliography entry of the article is transcribed and no reference is redistributed. Reference adaptation: none was needed and none was made; every reference inside the delivered unit resolves inside it, so no unresolved reference or citation is produced. Printed slips kept literal: no printed word, symbol, hypothesis or number of the counted statement or of its proof was altered. Layout and class: the article's own class is \texttt{amsart} with packages including amsfonts, amsthm, amsmath, amssymb, amscd, mathabx, esint, graphics, indentfirst, latexsym, epsfig, mathrsfs, color, enumitem; the wrapper is the lane's pinned \texttt{amsart} editorial wrapper at 10pt on A4, which restates the theorem-like environments the retained text needs and adds the article's own macro definitions that the retained text needs (\texttt{\textbackslash R}, \texttt{\textbackslash T}, \texttt{\textbackslash X}, \texttt{\textbackslash Y}, \texttt{\textbackslash aul}, \texttt{\textbackslash aur}, \texttt{\textbackslash avl}, \texttt{\textbackslash avr}, \texttt{\textbackslash norm}, \texttt{\textbackslash p}, \texttt{\textbackslash ph}, \texttt{\textbackslash supt}, \texttt{\textbackslash ut}, \texttt{\textbackslash vsx}, \texttt{\textbackslash vt}, \texttt{\textbackslash w}, \texttt{\textbackslash wh}), with extra wrapper packages (bm, bbm, dsfont, xspace, cancel, mathtools). The delivered numbering is the unit's own. Assets: no retained span contains a figure environment, a graphics-inclusion command, a PDF-page inclusion, a table or a TikZ picture; no image file is copied into this package, the package declares no asset dependency and no image file is redistributed with it. Licence: version arXiv:2507.21470v1 of this article is distributed under the Creative Commons Attribution 4.0 International licence, as recorded in the arXiv licence metadata for this exact version and shown on the abstract page https://arxiv.org/abs/2507.21470v1. A full rights notice is delivered at licenses/arxiv-2507.21470-CC-BY-4.0.txt. Characters: every retained excerpt is pure ASCII, so the delivered engine, XeLaTeX with its default Latin Modern text font, reports no missing character.
\begin{equation}\label{VNC-Ion}
	\begin{cases}
		\p_tv-\p_xu=0,&\quad\\
		\p_tu-\p_x{\sigma(v)}=\mu\p^2_xu,&\quad\\
	\end{cases}x\in \R, t>0,
\end{equation}

\begin{equation}\label{initial}
	(v,u)(x,0)=(v_0(x),u_0(x)),~~x\in\R,
\end{equation} 

\begin{equation}\label{cond-farfield}
	(v_0,u_0)(x)~\to 
	\begin{cases}
		(\overline{v}_l,\overline{u}_l)+(\phi_{0l},\psi_{0l}) \qquad \text{as}~x~\rightarrow -\infty\\
		(\overline{v}_r,\overline{u}_r)+(\phi_{0r},\psi_{0r}) \qquad \text{as}~x~\rightarrow +\infty
	\end{cases}
\end{equation}	

\begin{equation}\label{0-A}
	\int_{0}^{\pi_l}(\phi_{0l},\psi_{0l})\,dx=0 \quad and \quad \int_{0}^{\pi_r}(\phi_{0r},\psi_{0r})\,dx=0.
\end{equation}

\begin{lemma}\label{decay}
	Assume that $(v_0,u_0)(x) \in H^k(0,\pi)$ with $k\geqq 2$ is periodic with period $\pi > 0$ and average $(\overline{v},\overline{u})$. Then there
	exists $\varepsilon_0>0$ such that if 
	\begin{align*}
		\varepsilon:=\norm{(v_0,u_0)-(\overline{v},\overline{u})}_{H^k(0,\pi)}\leq \varepsilon_0,
	\end{align*}
	the problem \eqref{VNC-Ion} with initial data $(v_0,u_0)$ admits a unique periodic solution
	\begin{align*}
		(v,u)(x,t)\in C(0,+\infty;H^k(0,\pi)),
	\end{align*}
	which has some period and average as $(v_0,u_0)$. Moreover, it holds that 
	\begin{equation}\label{decay1}
		\norm{(v,u)-(\overline{v},\overline{u})}_{H^k(0,\pi)}(t)\leq C\varepsilon e^{-\alpha t},~t\geqq 0,
	\end{equation}
	where the constant $C>0$ and $\alpha>0$ are independent of $\varepsilon$ and t.
\end{lemma}

\begin{equation}\label{XY}
	\begin{cases}
		\begin{aligned}
			&\X'(t)=-s-\frac{\int_{\R}(u_r-u_l)g'_{st+\X}(x)  \,dx }{\int_{\R}(v_r-v_l)g'_{st+\X}(x)  \,dx },\\
			&\Y'(t)=-s-\frac{\int_{\R}[\sigma(v_r)-\sigma(v_l)+\mu\p_x u_r-\mu\p_x u_l]g'_{st+\Y}(x)\,dx }{\int_{\R}(u_r-u_l)g'_{st+\Y}(x) \,dx}.
		\end{aligned}
	\end{cases}
\end{equation}

\begin{lemma}\label{XYdecay}
	Assume that the periodic perturbations $\phi_{0l},\phi_{0r},\psi_{0l}$ and $\psi_{0r}$ satisfy \eqref{0-A}. Then there exists $\varepsilon_0>0$ such that
	if 
	\begin{align}\label{small assumption for shifts}
		\varepsilon:=\norm{\phi_{0l},\psi_{0l}}_{H^2(0,\pi_l)}+\norm{\phi_{0r},\psi_{0r}}_{H^2(0,\pi_r)}\leq \varepsilon_0,
	\end{align}
	there exists a unique solution $(\X,\Y)(t) \in C^1[0,+\infty)$ to \eqref{XY} with the initial data $(\X_0,\Y_0)$. Moreover, the solution satisfies that
	\begin{center}
		\begin{minipage}[c]{7cm}
			$ |\X'(t)|+\left|\mathcal{X}(t)-\mathcal{X}_{\infty}\right| \leq  C \varepsilon e^{-\alpha t},$ \\
		  $|\Y'(t)|+\left|\mathcal{Y}(t)-\mathcal{Y}_{\infty}\right| \leq C \varepsilon e^{-\alpha t},$
		\end{minipage}
		\begin{minipage}[c]{2cm} 
			$ t \geqq 0,$
		\end{minipage}
	\end{center}
	where the constants $C>0$ and $\alpha>0$ are independent of $\varepsilon$ and t, and the shifts $\X_{\infty}$ and $\Y_{\infty}$ are given by $\X_{\infty}=\mathcal{A}_1(\X_0)+\mathcal{C} _1$
	and $\Y_{\infty}=\mathcal{A} _2(\Y_0)+\mathcal{C} _2$, with
	\begin{equation}\label{C1}
		\mathcal{C} _1:=\frac{1}{\avr-\avl}\left\{\frac{1}{\pi_l}\int_{0}^{\pi_l} \int_{0}^{x} \phi_{0l} \,dy  \,dx -\frac{1}{\pi_r}\int_{0}^{\pi_r} \int_{0}^{x} \phi_{0r} \,dy  \,dx\right\},
	\end{equation}
	\begin{equation}\label{C2}
		\begin{aligned}
		&\mathcal{C}_2:=\frac{1}{\aur-\aul}\left\{\frac{1}{\pi_l}\int_{0}^{\pi_l} \int_{0}^{x} \psi_{0l}(y) \,dy  \,dx+\int_{0}^{+\infty}\frac{1}{\pi_l}\int_{0}^{\pi_l}[\sigma(v_l(x,t))-\sigma(\avl)] \,dx \,dt\right.\\ 
		&\left.-\frac{1}{\pi_r}\int_{0}^{\pi_r} \int_{0}^{x} \psi_{0r}(y) \,dy  \,dx-\int_{0}^{+\infty}\frac{1}{\pi_r}\int_{0}^{\pi_r}[\sigma(v_r(x,t))-\sigma(\avr)] \,dx \,dt\right\}
	    \end{aligned}
	\end{equation}
\end{lemma}

\begin{equation}\label{new0mass}
	\begin{aligned}
		&s\left\{\int_{-\infty}^{0}(v_0-v^S-\phi_{0l})(x)\,dx+\int_{0}^{+\infty} (v_0-v^S-\phi_{0r})(x)\,dx\right.\\
		&\left.-\frac{1}{\pi_l}\int_{0}^{\pi_l}\int_{0}^{x} \phi_{0l}(y) \,dy \,dx+\frac{1}{\pi_r}\int_{0}^{\pi_r}\int_{0}^{x} \phi_{0r}(y) \,dy \,dx\right\}\\
		&=-\int_{-\infty}^{0}(u_0-u^S-\psi_{0l})(x)\,dx-\int_{0}^{+\infty}(u_0-u^S-\psi_{0r})(x)\,dx\\
		&+\frac{1}{\pi_l}\int_{0}^{\pi_l}\int_{0}^{x} \psi_{0l}(y) \,dy \,dx+\int_{0}^{+\infty}\frac{1}{\pi_l}\int_{0}^{\pi_l}[\sigma(v_l(x,t))-\sigma(\avl)] \,dx \,dt\\
		&-\frac{1}{\pi_r}\int_{0}^{\pi_r}\int_{0}^{x} \psi_{0r}(y) \,dy \,dx-\int_{0}^{+\infty}\frac{1}{\pi_r}\int_{0}^{\pi_r}[\sigma(v_r(x,t))-\sigma(\avr)] \,dx \,dt.
	\end{aligned}
\end{equation}

\begin{theorem}\label{mainresult}
	Assume that the periodic perturbations in \eqref{cond-farfield} satisfy \eqref{0-A} and \eqref{new0mass}. Then there exists $\varepsilon_0>0$ such that
	, if $E_0\leq \varepsilon_0$, there admits a unique global solution of \eqref{VNC-Ion} and \eqref{initial}, satisfying
\begin{equation}\label{regularity}
		\begin{aligned}
		v-\vt \in C\Big([0,+\infty);H^1(\R)\Big)\cap L^2\Big((0,+\infty);H^1(\R)\Big),\\
		u-\ut \in C\Big([0,+\infty);H^1(\R)\Big)\cap L^2\Big((0,+\infty);H^2(\R)\Big),
	\end{aligned}
\end{equation}
	and 
	\begin{equation}\label{asymptotic}
		\norm{(v,u)(\cdot,t)-(v^S,u^S)(\cdot-st-\X_{\infty})}_{L^{\infty}(\R)}\to 0 \quad \text{as}~t \to +\infty,
	\end{equation}
	with the constant shift $\X_{\infty}(=\Y_{\infty})$ given by
	\begin{align*}
		\X_{\infty}=-\frac{1}{\avr-\avl}\left\{\int_{-\infty}^{0}(v_0-v^S-\phi_{0l})(x)  \,dx+\int_{0}^{+\infty}(v_0-v^S-\phi_{0r})(x)  \,dx  \right\}+\mathcal{C}_1,
	\end{align*}
	where $\mathcal{C} _1$ is given in \eqref{C1}.
\end{theorem}

\begin{equation}\label{system2}
	\begin{cases}
		\p_t\Phi-s\p_{\xi}\Phi-\p_{\xi}\Psi=\mathbf{H}_1,\\
		\p_t\Psi-s\p_{\xi}\Psi-\sigma'(\mathbf{\tilde{v}})\p_{\xi}\Phi=\sigma(\mathbf{v}|\mathbf{\tilde{v}})+\mu \p^2_{\xi}\Psi+\mathbf{H}_2,
	\end{cases}
\end{equation}

\begin{equation}\label{hterm}
	\sigma(\mathbf{v}|\mathbf{\tilde{v}}):=\sigma(\mathbf{v})-\sigma(\mathbf{\tilde{v}})-\sigma'(\mathbf{\tilde{v}})\p_{\xi}\Phi\geqq 0
\end{equation}

\begin{equation}\label{newinitial}
	(\Phi,\Psi)(\xi,0)=(\Phi_0,\Psi_0)(\xi) \in H^2(\R) \times H^2(\R),
\end{equation}

\begin{proposition}\label{mainpro}
	(A priori estimates). For any $\T>0$, assume that $(\Phi,\Psi)$ is a smooth solution to \eqref{system2} and \eqref{newinitial}. Then there
	exist $\varepsilon_0>0$ and $\delta_0>0$, independent of ~$\T$, such that if 
	\begin{equation}\label{initialsmall}
		\varepsilon=E_0<\varepsilon_0~~\text{and}~~\delta=\mathop{sup}\limits_{t \in (0,\T)}\norm{(\Phi,\Psi)}_2(t)<\delta_0,
	\end{equation}
	then
	\begin{equation}\label{proesti}
		\mathop{sup}\limits_{t \in [0,\T]}\norm{(\Phi,\Psi)}^2_2(t)+\int_{0}^{\T}(\norm{\p_{\xi}\Phi}^2_1+\norm{\p_{\xi}\Psi}^2_2)(t)  \,dt\leq C_0
		(\norm{(\Phi_0,\Psi_0)}^2_2+\varepsilon),
	\end{equation}
	where $C_0>0$ is independent of $\varepsilon$,$\delta$ and $\T$.
\end{proposition}

\begin{equation}\label{small1}
	\begin{aligned}
		\mathop{sup}\limits_{t \in [0,\T]}\norm{(\Phi,\Psi)}_{W^{1,\infty}(\R)}\leq C\mathop{sup}\limits_{t \in [0,\T]}\norm{(\Phi,\Psi)}_2\leq C\delta_0,\\
		\mathop{sup}\limits_{t \in [0,\T]}\norm{W}_{L^{\infty}(\R)}\leq C\mathop{sup}\limits_{t \in [0,\T]}\norm{W}_1\leq C\delta_0,
	\end{aligned}
\end{equation}

\begin{lemma}\label{decay-error}
	Under the assumptions of Theorem \ref{mainresult}, it holds that
	\begin{align*}
		\norm{\mathbf{H}_1}_2(t),\norm{\mathbf{H}_2}_1(t)\leq C\varepsilon e^{-\alpha t},~t\geqq 0,
	\end{align*}
	where $C>0$ is independent of $\varepsilon$ and $t$, and $\alpha>0$ is the constant given in Lemma \ref{XYdecay}.
\end{lemma}

  \begin{equation}\label{XYZ-cond}
  	X(v^S)\geqq 0, ~~Y(v^S)\geqq0,~~Z(v^S)\geqq C|\sigma''(v^S)|\text{ in }[\avl,\avr].
  \end{equation}

\begin{equation}\label{transe}
	\begin{aligned}
		&\p_t\left(\frac{a(\vsx)}{2}\ph^2+\frac{w(\vsx)}{2}\w^2\right)+\mu a(\vsx)(\p_{\xi}\ph)^2\\
		&+\mu X(\vsx)\ph^2+\frac{a(\vsx)}{s}(\vsx)'Y(\vsx)(s\ph+\w)^2+(\vsx)'Z(\vsx)\w^2\\
		&=\underbrace{\frac{\X'(\vsx)'(\ph^2(2a(\vsx)T'(\vsx)-a'(\vsx)T(\vsx))+\w^2(2w(\vsx)T'(\vsx)-w'(\vsx)T(\vsx)))}{2T(\vsx)}}_{\mathbf{I}_{1,1}}\\
		&+\underbrace{\frac{\sigma(\mathbf{v}|\mathbf{\tilde{v}})w(\vsx)\w+(\mathbf{H}_2-\mu\p_{\xi}\mathbf{H}_1)w(\vsx)\w}{T(\vsx)}}_{\mathbf{I}_{1,2}}+\underbrace{\frac{\mathbf{H}_1\ph a(\vsx)}{T(\vsx)}}_{\mathbf{I}_{1,3}}+\underbrace{(\sigma'(\mathbf{\tilde{v}})-\sigma'(\vsx))w(\vsx)\p_{\xi}\ph\w}_{\mathbf{I}_{1,4}}\\
		&+\underbrace{\frac{T'(\vsx)}{T(\vsx)}(\sigma'(\mathbf{\tilde{v}})-\sigma'(\vsx))w(\vsx)(\vsx)'\ph\w}_{\mathbf{I}_{1,5}}+\underbrace{\p_{\xi}\zeta(\xi)}_{\mathbf{I}_{1,6}},\\
	\end{aligned}
\end{equation}

\begin{equation}\label{Hbig}
	H\geq\frac{4(s^2+\sigma'(v_b))\sigma''(v_b)}{5(s^2-\sigma'(v_b))},
\end{equation}

\begin{lemma}\label{construction}
	Under the assumption of Proposition \ref{mainpro}, there exists $\varepsilon_0>0$ and $\delta_0>0$ and such that if $\varepsilon<\varepsilon_0$ and $\delta<\delta_0$, then we have the estimate as follows,
	\begin{equation}\label{transformestimate}
		\begin{aligned}
			&\mathop{sup}\limits_{t \in [0,\T]}\norm{\Phi,W}^2(t) +\int_{0}^{\T}\int_{\R} |(\vsx)'\sigma''(\vsx)|W^2\,d\xi \,dt\\
			&\leq C(\varepsilon+\norm{\Phi_0,W_0}^2+\delta\int_{0}^{T} \norm{\p_{\xi}\Phi}^2 \,dt ),
		\end{aligned}
	\end{equation}
	where $C>0$ is independent of $\varepsilon,\delta$ and $\T$.
\end{lemma}

\begin{proof}
	The estimate \eqref{transformestimate} would follows from integrating \eqref{transe} with respect to $\xi$ on $\mathbb{R}$, \eqref{XYZ-cond} and estimates for $\textbf{I}_{1,i}$,$i=1,\cdots,6$. Due to our construction of the weighted and transform functions, we will integrate separately with respect to $\xi$ on the intervals for $\xi$ corresponding to $\vsx$ in $[0,\avr]$ and $[\avl,0]$. Since $\vsx$ is an increasing function of $\xi$, let
	$$\vsx(\xi_i+\X(t))=v^S(\xi_i)=v_i, i=-1,0,1,\cdots,N$$
	$$\vsx(\xi_d+\X(t))=v^S(\xi_d)=v_d$$
	with $v_i=\frac{i}{2H},i=-1,\cdots, N-1$, $v_N=v_b$, $v_d$ defined in section 3.1.2.  
	
	First, we integrate \eqref{transe} over $[0,\T]\times[\xi_{-1}+\X(t),\xi_0+\X(t))$,  \dots, $[0,\T]\times[\xi_{N-1}+\X(t),\xi_N+\X(t)]$ and 
	$[0,\T]\times[\xi_N+\X(t),+\infty)$ and add the resulting equations together to get
\begin{equation}\label{pst-part}
	\begin{aligned}
		&\int_{\xi_{-1}+\X(t)}^{+\infty}\left(\frac{a(\vsx)}{2}\ph^2+\frac{w(\vsx)}{2}\w^2\right)(t)\,d\xi+\int_{0}^{\T}\int_{\xi_{-1}+\X(t)}^{+\infty}\mu a(\vsx)(\p_{\xi}\ph)^2\,d\xi dt\\
		&+\int_{0}^{\T}\int_{\xi_{-1}+\X(t)}^{+\infty}\mu X(\vsx)\ph^2+\frac{a(\vsx)}{s}(\vsx)'Y(\vsx)(s\ph+\w)^2+(\vsx)'Z(\vsx)\w^2\,d\xi dt\\
		&=\sum_{i=1}^{6}\int_{0}^{\T}\int_{\xi_{-1}+\X(t)}^{+\infty}\textbf{I}_{1,i}\,d\xi dt.
	\end{aligned}
\end{equation}
	It remains to control the right hand side of \eqref{pst-part}. By the construction of $T(\vsx),w(\vsx)$ for $\vsx\in[v_{-1},\avr]$, we have $T(\vsx),w(\vsx) \in C[v_{-1},\avr]$,
	but not in $C^1[v_{-1},\avr]$.  Thus,
	$$\zeta(\xi_{-1}+\X(t))=\zeta(+\infty)=0,$$
	and
	$$[\zeta](\xi_i+\X(t)):=\zeta_{-\X}(\xi_i-)-\zeta_{-\X}(\xi_i+)=\mu aY(\vsx)'\ph^2(v_i^+)-\mu aY(\vsx)'\ph^2(v_i^-),i=1,\cdots, N.$$
	It is easy to check (c.f. Appendix) that
	\begin{equation}\label{interact}
		\begin{aligned}
			&[\zeta](v_i)=0,~~i=1,2,3,\dots,N-1,\\
			&[\zeta](v_N)=\frac{\mu}{2}[-\frac{5H}{2^{\frac{13}{4}}}+\frac{(s^2+\sigma'(v_b))\sigma''(v_b)}{2^{\frac{5}{4}}(s^2-\sigma'(v_b))}]\leq 0,
		\end{aligned}
	\end{equation}
	for $H$ being chosen large enough in section 3.1.2, especially in \eqref{Hbig}. 
 Therefore, 
	\begin{equation}\label{flux-term1}
		\int_{0}^{\T}\int_{\xi_{-1}+\X(t)}^{+\infty}\textbf{I}_{1,6}\,d\xi dt=\int_{0}^{\T}\sum_{i=1}^{N}[\zeta](v_i)\leq 0
	\end{equation}
For other terms $\int_{0}^{\T}\int_{\xi_{-1}+\X(t)}^{+\infty}\textbf{I}_{1,i}\,d\xi dt,i=1,\cdots,5$,	it follows from Lemma \ref{decay} and \ref{XYdecay} that
	\begin{equation}\label{traveling-wave-decay}
		\norm{\vsx-\mathbf{\tilde{v}}}_{L^{\infty}(\R)}=\norm{(\mathbf{v}_l-\avl)(1-g_{\X})+(\mathbf{v}_r-\avr)g_{\X}}_{L^{\infty}(\R)}\leq C\varepsilon e^{-\alpha t}.
	\end{equation}
	\begin{equation}\label{X'decay}
		|\X'(t)|\leq C\varepsilon e^{-\alpha t}.
	\end{equation}
Note that $T(\vsx),w(\vsx),T'(\vsx),w'(\vsx)$ is bounded for $\vsx\in[v_{i},v_{i+1}]$, $i=-1,\cdots,N$, and$\vsx\in[v_N, \avr]$, and $T(\vsx)>0$ is bounded from below. By \eqref{X'decay}, we have 
\begin{equation}\label{X'est}
	\begin{aligned}
		\int_{0}^{\T}\int_{\xi_{-1}+\X(t)}^{+\infty}|\mathbf{I}_{1,1}|\,d\xi\,dt&\leq \int_{0}^{\T}\int_{\xi_{-1}+\X(t)}^{+\infty}|\X'(t)|(\ph^2+\w^2)\,d\xi\,dt\\
		&\leq C\varepsilon\int_{0}^{\T}e^{-\alpha t}(\norm{\ph}^2+\norm{\w}^2) \,dt \\
		&\leq C\varepsilon(\supt\norm{\ph}^2+\supt\norm{\w}^2).
	\end{aligned}
\end{equation}
	By \eqref{hterm}, one has $|\sigma(\mathbf{v}|\mathbf{\tilde{v}})|\leq C|\p_{\xi}\Phi|^2$. Then, we can get from \eqref{small1} and Lemma  \ref{decay-error} that,
	\begin{equation}\label{nondegenerate1}
		\begin{aligned}
			&\int_{0}^{\T} \int_{\xi_{-1}+\X(t)}^{+\infty}|\mathbf{I}_{1,2}|\,d\xi\,dt\leq\int_{0}^{\T} \int_{\xi_{-1}+\X(t)}^{+\infty}|(|\p_\xi\Phi|^2+|\mathbf{H}_2|+|\p_{\xi}\mathbf{H}_1|)|\w|\,d\xi\,dt\\
			&\leq C\supt\norm{W}_{L^{\infty}(\R)}\int_{0}^{\T}\norm{\p_{\xi}\Phi}^2\,dt+C\int_{0}^{\T} (\norm{\p_{\xi}\mathbf{H}_1}+\norm{\mathbf{H}_2})\norm{\w} \,dt\\
			&\leq C\delta\int_{0}^{\T}\norm{\p_{\xi}\Phi}^2\,dt+C\varepsilon\int_{0}^{\T} e^{-\alpha t}\norm{\w} \,dt\\
			&\leq C\delta\int_{0}^{\T}\norm{\p_{\xi}\Phi}^2\,dt+C\varepsilon+C\varepsilon \supt\norm{\w}^2,
		\end{aligned}
	\end{equation}
	and
	\begin{equation}\label{nondegenerate2}
		\begin{aligned}
			\int_{0}^{\T} \int_{\xi_{-1}+\X(t)}^{+\infty}|\mathbf{I}_{1,3}|\,d\xi\,dt\leq C\int_{0}^{\T}\norm{\mathbf{H}_1}\norm{\ph}  \,dt\leq C\varepsilon+C\varepsilon \supt\norm{\ph}^2.
		\end{aligned}
	\end{equation}
	By \eqref{traveling-wave-decay}, we have
	\begin{equation}\label{nondegenerate3}
		\begin{aligned}
			\int_{0}^{\T} \int_{\xi_{-1}+\X(t)}^{+\infty}|\mathbf{I}_{1,4}|\,d\xi\,dt&\leq C\int_{0}^{\T} \int_{\xi_{-1}+\X(t)}^{+\infty}|\mathbf{\tilde{v}}-\vsx||\p_{\xi}\ph||\w|\,d\xi\,dt\\
			&\leq C\varepsilon\int_{0}^{\T}e^{-\alpha t}(\norm{\p_{\xi}\ph}^2+\norm{\w}^2)  \,dt\\
			&\leq C\varepsilon\int_{0}^{\T} \norm{\p_{\xi}\ph}^2 \,dt+C\varepsilon\supt\norm{\w}^2.
		\end{aligned}
	\end{equation}
	Similarly, by \eqref{traveling-wave-decay} and \eqref{X'decay}, one has that
	\begin{equation}\label{nondegenerate4}
		\begin{aligned}
			\int_{0}^{\T} \int_{\xi_{-1}+\X(t)}^{+\infty}|\mathbf{I}_{1,5}|\,d\xi\,dt&\leq C\int_{0}^{\T} \int_{\xi_{-1}+\X(t)}^{+\infty}|\mathbf{\tilde{v}}-\vsx||\ph||\w|\,d\xi\,dt\\
			&\leq C\varepsilon\int_{0}^{\T} e^{-\alpha t}(\norm{\ph}^2+\norm{\w}^2) \,dt\\
			&\leq C\varepsilon(\supt\norm{\ph}^2+\supt\norm{\w}^2).
		\end{aligned}
	\end{equation}
	Plugging \eqref{flux-term1},\eqref{X'est}-\eqref{nondegenerate4} into \eqref{pst-part} and using \eqref{XYZ-cond} for $\vsx\in[v_{-1},\avr]$, we have
	\begin{equation}\label{pst-integral}
		\begin{aligned}
			&\int_{\xi_{0}+\X(t)}^{+\infty}\left(\frac{a(\vsx)}{2}\ph^2+\frac{w(\vsx)}{2}\w^2\right)(t)\,d\xi+\int_{0}^{\T}\int_{\xi_{0}+\X(t)}^{+\infty}(\vsx)'|\sigma''(\vsx)|\w^2\,d\xi dt\\
			&\leq\int_{\xi_{-1}+\X(t)}^{+\infty}\left(\frac{a(\vsx)}{2}\ph^2+\frac{w(\vsx)}{2}\w^2\right)(t)\,d\xi+\int_{0}^{\T}\int_{\xi_{-1}+\X(t)}^{+\infty}(\vsx)'|\sigma''(\vsx)|\w^2\,d\xi dt\\
			&\leq C\varepsilon(\supt\norm{\ph}^2+\supt\norm{\w}^2)+C\varepsilon+C\delta\int_{0}^{\T}\norm{\p_{\xi}\Phi}^2\,dt
		\end{aligned}
	\end{equation}
	where the first inequality follows from the non-negative of integrands.
	
	 Next, we integrate \eqref{transe} $[0,\T]\times[-\infty,\xi_0+\X(t))$ and $[0,\T]\times[\xi_0+\X(t),\xi_d+\X(t))$, and add the resulting equations together to yield that
	 \begin{equation}\label{neg-part}
	 	\begin{aligned}
	 		&\int_{-\infty}^{\xi_d+\X(t)}\left(\frac{\hat{a}(\vsx)}{2}\ph^2+\frac{\hat{w}(\vsx)}{2}\w^2\right)(t)\,d\xi+\int_{0}^{\T}\int_{\xi_{-1}+\X(t)}^{+\infty}\mu \hat{a}(\vsx)(\p_{\xi}\ph)^2\,d\xi dt\\
	 		&+\int_{0}^{\T}\int_{-\infty}^{\xi_d+\X(t)}\mu X(\vsx)\ph^2+\frac{a(\vsx)}{s}(\vsx)'Y(\vsx)(s\ph+\w)^2+(\vsx)'Z(\vsx)\w^2\,d\xi dt\\
	 		&=\sum_{i=1}^{6}\int_{0}^{\T}\int_{-\infty}^{\xi_d+\X(t)}\textbf{I}_{1,i}\,d\xi dt.
	 	\end{aligned}
	 \end{equation}
	
%	Then we estimate the rest part, set $w(\vsx)=\wh(\vsx)$, $T(\vsx)=\th(\vsx)$, then $(\Phi,W)(\xi,t)=\th(\vsx(\xi))(\phh(\xi,t),\whh(\xi,t))$, substisubstituting these into \eqref{transe}, we have
%	\begin{equation}\label{transe1}
%		\begin{aligned}
%			&\p_t(\frac{\ah(\vsx)}{2}\phh^2+\frac{\wh(\vsx)}{2}\whh^2)+\mu(\p_{\xi}\phh)^2\ah(\vsx)+\whh^2(\vsx)'Z(\vsx)\\
%			&+Y(\vsx)(\vsx)'(s\phh+\whh)^2+\mu\phh^2X(\vsx)\\
%			&=\underbrace{\frac{\X'(\vsx)'(\phh^2(2\ah(\vsx)\th'(\vsx)-\ah'(\vsx))+\whh^2(2\wh(\vsx)\th'(\vsx)-\wh'(\vsx)))}{2\th(\vsx)}}_{\mathbf{I}_{1,7}}\\
%			&+\underbrace{\frac{\sigma(\mathbf{v}|\mathbf{\tilde{v}})\wh(\vsx)\whh+(\mathbf{H}_2-\mu\p_{\xi}\mathbf{H}_1)\wh(\vsx)\whh}{\th(\vsx)}}_{\mathbf{I}_{1,8}}+\underbrace{\frac{\mathbf{H}_1\phh \ah(\vsx)}{\th(\vsx)}}_{\mathbf{I}_{1,9}}+\underbrace{(\sigma'(\mathbf{\tilde{v}})-\sigma'(\vsx))\wh(\vsx)\p_{\xi}\phh\whh}_{\mathbf{I}_{1,10}}\\
%			&+\underbrace{\frac{\th'(\vsx)}{\th(\vsx)}(\sigma'(\mathbf{\tilde{v}})-\sigma'(\vsx))\wh(\vsx)(\vsx)'\phh\whh}_{\mathbf{I}_{1,11}}+\underbrace{\p_{\xi}\left\{\dots\right\}}_{\mathbf{I_{1,12}}},\\
%		\end{aligned}
%	\end{equation}
%	where $\left\{\dots\right\}=\frac{s \ah(\vsx)}{2}\phh^2+\frac{s\wh(\vsx)}{2}\whh^2+\whh\phh \ah(\vsx)+\mu\frac{\th'(\vsx)}{\th(\vsx)}\ah(\vsx)(\vsx)'\phh^2-\frac{\mu}{2}\ah'(\vsx)(\vsx)'\phh^2+\mu \ah(\vsx)\p_{\xi}\phh\phh$.
	Note that $T(\vsx),w(\vsx) \in C^1[\avl,v_d]$, it is easy to check (c.f. Appendix) that
	\begin{equation}\label{interact2}
		\zeta(-\infty)=0,\quad \zeta(\xi_d+\X(t))\leq 0.
	\end{equation}
	Therefore, 
	\begin{equation}\label{flux-term2}
		\int_{0}^{\T}\int_{-\infty}^{\xi_{d}+\X(t)}\textbf{I}_{1,6}\,d\xi dt\leq 0.
	\end{equation}
%	 Integrate $\mathbf{I_{1,12}}$ over $[0,\T] \times (-\infty,\xi_0+\X(t)]$ and $[0,\T] \times [\xi_0+\X(t),\xi_d+\X(t)]$ we find it is non-positive, the details are also left in "Appendix".
%	
%	Set $\vsx(\xi_d+\X(t))=v_d$, we find $\wh(\vsx),\th(\vsx) \in C^1[\avl,v_d]$ but $\th''(\vsx)$ is not continuous at the point $\xi_0+\X(t)$, and we have
%\begin{align}
%	&X(\vsx):=\p_{\xi}\th\p_{\xi}\left(\frac{\ah}{\th}\right)-\frac{1}{2}\p_{\xi}^2\ah\nonumber\\
%	&=\big[\frac{\th'(\vsx)}{\th(\vsx)}\ah'(\vsx)-\big(\frac{\th'(\vsx)}{\th(\vsx)}\big)^2\ah(\vsx)-\frac{\ah''(\vsx)}{2}\big]((\vsx)')^2-\frac{\ah'(\vsx)(\vsx)''}{2}\nonumber\\
%	&=\big[\frac{\th'(\vsx)}{\th(\vsx)}\ah'(\vsx)-\big(\frac{\th'(\vsx)}{\th(\vsx)}\big)^2\ah(\vsx)-\frac{\ah''(\vsx)}{2}\big](\frac{h(\vsx)}{s\mu})^2-\frac{\ah'(\vsx)h(\vsx)h'(\vsx)}{2(s\mu)^2},
%\end{align}
%\begin{equation}
%	Y(\vsx):=\frac{\ah'(\vsx)}{2\ah(\vsx)}-\frac{\th'(\vsx)}{\th(\vsx)},
%\end{equation}
%\begin{align}
%	Z(\vsx):=\frac{s\wh'(\vsx)}{2}-s\frac{\th'(\vsx)}{\th(\vsx)}w(\vsx)+\frac{\th'(\vsx)}{s\th(\vsx)}\ah(\vsx)-\frac{\ah'(\vsx)}{2s}.
%\end{align}

Similar to \eqref{X'est}-\eqref{nondegenerate4}, we have the following estimates for other terms on the right hand side of \eqref{neg-part}
	\begin{equation}\label{hX'est}
		\begin{aligned}
			\int_{0}^{\T}\int_{-\infty}^{\xi_{d}+\X(t)}|\mathbf{I}_{1,1}|\,d\xi\,dt\leq C\varepsilon(\supt\norm{\ph}^2+\supt\norm{\w}^2),
		\end{aligned}
	\end{equation}
	\begin{equation}\label{hnondegenerate1}
			\int_{0}^{\T} \int_{-\infty}^{\xi_{d}+\X(t)}|\mathbf{I}_{1,2}|\,d\xi\,dt\leq C\delta\int_{0}^{\T}\norm{\p_{\xi}\Phi}^2\,dt+C\varepsilon+C\varepsilon \supt\norm{\w}^2,
	\end{equation}
	\begin{equation}\label{hnondegenerate2}
		\begin{aligned}
			\int_{0}^{\T} \int_{-\infty}^{\xi_{d}+\X(t)}|\mathbf{I}_{1,3}|\,d\xi\,dt\leq C\varepsilon+C\varepsilon \supt\norm{\ph}^2,
		\end{aligned}
	\end{equation}
	\begin{equation}\label{hnondegenerate3}
		\begin{aligned}
			\int_{0}^{\T} \int_{-\infty}^{\xi_{d}+\X(t)}|\mathbf{I}_{1,4}|\,d\xi\,dt
			\leq C\varepsilon\int_{0}^{\T} \norm{\p_{\xi}\ph}^2 \,dt+C\varepsilon\supt\norm{\w}^2,
		\end{aligned}
	\end{equation}
	\begin{equation}\label{hnondegenerate4}
		\begin{aligned}
			\int_{0}^{\T} \int_{-\infty}^{\xi_{d}+\X(t)}|\mathbf{I}_{1,5}|\,d\xi\,dt\leq C\varepsilon(\supt\norm{\ph}^2+\supt\norm{\w}^2).
		\end{aligned}
	\end{equation}
	Plugging \eqref{flux-term2}-\eqref{hnondegenerate4} into \eqref{neg-part} and using \eqref{XYZ-cond} for $\vsx\in[\avl,v_d]$, we have
	\begin{equation}\label{neg-integral}
		\begin{aligned}
			&\int_{-\infty}^{\xi_{0}+\X(t)}\left(\frac{\hat{a}(\vsx)}{2}\ph^2+\frac{\hat{w}(\vsx)}{2}\w^2\right)(t)\,d\xi+\int_{-\infty}^{\xi_{d}+\X(t)}(\vsx)'|\sigma''(\vsx)|\w^2\,d\xi dt\\
			&\leq\int_{-\infty}^{\xi_{0}+\X(t)}\left(\frac{\hat{a}(\vsx)}{2}\ph^2+\frac{\hat{w}(\vsx)}{2}\w^2\right)(t)\,d\xi+\int_{-\infty}^{\xi_{d}+\X(t)}(\vsx)'|\sigma''(\vsx)|\w^2\,d\xi dt\\
			&\leq C\varepsilon(\supt\norm{\ph}^2+\supt\norm{\w}^2)+C\varepsilon+C\delta\int_{0}^{\T}\norm{\p_{\xi}\Phi}^2\,dt
		\end{aligned}
	\end{equation}
	Note that $w(\vsx)$ and $\wh(\vsx)$ have uniform lower bound $min\big\{\frac{1}{2H},\frac{(v_d)^2}{\sigma'(0)}\big\}$ whenever $\vsx \in [0,\avr]$ and $\vsx \in [\avl,0]$, and $T$ is also bounded from below and above so that $\norm{\ph,\w} \sim \norm{\Phi,W}$.  Summing \eqref{pst-integral} with \eqref{neg-integral} and using the lower bounds of $w$ and $a$, we can get \eqref{transformestimate}.
%	
%	
%	If $\varepsilon_0$ is small enough, integrate \eqref{transe} over $[0,\T]\times[\xi_{-1}+\X(t),\xi_0+\X(t)]$, $[0,\T]\times[\xi_0+\X(t),\xi_1+\X(t)]$, \dots, $[0,\T]\times[\xi_{N-1}+\X(t),\xi_N+\X(t)]$ and 
%	$[0,\T]\times[\xi_N+\X(t),+\infty)$, integrate \eqref{transe1} over $[0,\T] \times (-\infty,\xi_0+\X(t)]$ and $[0,\T] \times [\xi_0+\X(t),\xi_d+\X(t)]$, then adding them, consider \eqref{interact}, and $\norm{\ph,\w} \sim \norm{\Phi,W} \sim \norm{\phh,\whh}$ by the boundedness of $T(\vsx)$ and $\th(\vsx)$, 
%	$w(\vsx)\sim T(\vsx) \sim \wh(\vsx) \sim \th(\vsx) \sim constant$.
%	
%	We also consider that when $\vsx \in [0,+\infty)$ and $\vsx \in (-\infty,0]$, $w(\vsx)$ and $\wh(\vsx)$ have uniform lower bound $min\big\{\frac{1}{2H},\frac{(v_d)^2}{\sigma'(0)}\big\}$, collecting the estimates from
%	$\mathbf{I}_{1,1}$ to $\mathbf{I}_{1,12}$, we finally have \eqref{transformestimate},
%	where $C$ is independent of $\varepsilon$, $\delta$ and $\T$.
\end{proof}

\end{document}
